\documentclass[../midgard.tex]{subfiles}
\graphicspath{{\subfix{../images/}}}
\begin{document}

\section{Witness Staking Script}
\label{h:witness}

A generic staking script on L1, parameterized by the unique identifier of a user event, whose registration state can be used to prove a falsely stated user event has not happened.

There are 3 logical endpoints this script needs to support:
\begin{enumerate}
    \item Validating registration for an authentic user event.
    \item Validating unregistration for the conclusion of a previously registered user event.
    \item Allowing an isolated registration that is immediately followed by an unregistration in the same transaction.\footnote{The witness uses the ordered certificate list; a certificate index is not a redeemer-list index.}
      Essentially proving the unregistered state of the staking script, without changing its state.
\end{enumerate}

\subsection{Staking script}
\label{h:witness-staking-script}
The \code{witness} script is parameterized by the unique ID of a corresponding user event (deposit, withdrawal, or transaction order).
Let \code{event\_id} be this identifier, which is the Blake2b256 hash of an L1 nonce output-reference spent in the user event's first transaction.

\begin{description}
  \item[Mint or Burn.] Script invocation to be included in the first transaction of a user event.
    Conditions:
      \begin{enumerate}
          \item Let \code{target\_policy} be the policy ID of the user event NFT, accessed via the redeemer.
          \item Registration/unregistration of this witness script is respectively tied to the \code{target\_policy}'s mint/burn logic. 
              \begin{itemize}
                  \item Registration is allowed if a user event NFT with asset name of \code{event\_id} is minted.
                  \item Unregistration is allowed if a user event NFT with asset name of \code{event\_id} is burnt.
              \end{itemize}
      \end{enumerate}
  \item[Register to Prove Not Registered.] The current redeemer carries \code{registration\_certificate\_index}.
    The certificate at this index must be the executing registration certificate.
    The immediately following certificate must unregister the same credential.
  \item[Unregister to Prove Not Registered.] The executing certificate must unregister a credential.
    The certificate at \code{registration\_certificate\_index} must register that same credential.
\end{description}

\paragraph{Conformance boundary.}
The current validator reads \code{transaction.certificates}; it does not enforce the former draft's additional reciprocal redeemer-data checks.
If those stronger checks are retained as a protocol requirement, they remain an implementation-conformance obligation rather than a current guarantee.
An unregistered credential proves absence of a live event witness at that ledger point, not that the event never existed: normal settlement also deregisters witnesses.
Fabricated-event proofs must bind their other live/confirmed-history evidence accordingly.

\end{document}
